\documentclass[11pt,aspectratio=169]{beamer}
\usepackage[utf8]{inputenc}
\usetheme{Madrid}
%%%%%%%%%%%%%%%%%%%%%%%%%%%%%%%%%%%%%%%%%%%%%%%%%%%%%%%%%%%%%%%%%%%%%%%%%%%%%%%%%%%%%%%%%%%%%%%%%%%%%%%%%%%%%%%%%%%%%%%%%%%%%%%%%%%%%%%
\definecolor{blue_text}{RGB}{0,115,197}
\definecolor{blue1}{RGB}{26,129,203}
\definecolor{blue2}{RGB}{102,171,220}
\definecolor{blue3}{RGB}{154,199,232}

\definecolor{green1}{RGB}{0,128,0}
\definecolor{green2}{RGB}{204,229,204}
\definecolor{green3}{RGB}{229,242,229}

\definecolor{lightblue2}{RGB}{204,227,243}
\definecolor{lightblue3}{RGB}{229,241,249}

\definecolor{orange0}{RGB}{255,116,0}
\definecolor{blue0}{RGB}{0,103,169}%vaidya
\definecolor{green0}{RGB}{0,148,104}


\setbeamercolor*{palette tertiary}{bg=blue1,fg=white}
\setbeamercolor*{palette secondary}{bg=blue2,fg=white}
\setbeamercolor*{palette primary}{bg=blue3,fg=black}

\setbeamercolor{titlelike}{bg=white,fg=blue_text}
\setbeamertemplate{title page}[default][rounded=false, shadow=false]

\setbeamercolor{frametitle}{bg=white,fg=blue_text}

\setbeamercolor{item projected}{bg=blue_text}

\setbeamertemplate{blocks}[rounded]%[shadow]
\setbeamercolor{block title example}{bg=green2,fg=green1}
\setbeamercolor{block body example}{bg=green3,fg=black}

\setbeamercolor{block title}{bg=lightblue2,fg=blue_text}
\setbeamercolor{block body}{bg=lightblue3,fg=black}

\setbeamercolor{button}{bg=blue2,fg=white}

\setbeamertemplate{navigation symbols}{}
%%%%%%%%%%%%%%%%%%%%%%%%%%%%%%%%%%%%%%%%%%%%%%%%%%%%%%%%%%%%%%%%%%%%%%%%%%%%%%%%%%%%%%%%%%%%%%%%%%%%%%%%%%%%%%%%%%%%%%%%%%%%%%%%%%%%%%%

\usepackage{eqnarray,amsmath}
\usepackage{amsfonts}
\usepackage{amssymb}
\usepackage{graphicx}
\usepackage{lmodern} 
\usepackage{bm}
\usepackage{epstopdf}
\usepackage{changepage}
\usepackage{array,booktabs}

\usepackage{pgfplots}
\usepackage{tikz}
\usetikzlibrary{patterns}
\usetikzlibrary{plotmarks}
\usepgfplotslibrary{fillbetween}

\DeclareMathOperator*{\argmax}{arg\,max}
\newcommand*\diff{\mathop{}\!\mathrm{d}}
\DeclareMathOperator{\supp}{supp}
\DeclareMathOperator{\prob}{Pr}

\usepackage{transparent}
\newcommand{\semitransp}[2][35]{\color{fg!#1}#2}

\usetikzlibrary{arrows.meta}

%%%%%%%%%%%%%%%%%%%%%%%%%%%%%%%%%%%%%%%%%%%%%%%%%%%%%%%%%%%%%%%%%%%%%%%%%%%%%%%%%%%%%%%%%%%%%%%%%%%%%%%%%%%%%%%%%%%%%%%%%%%%%%%%%%%%%%%


%====================================================
%========== DEFINITION OF AUTHORS ETC...=============
%====================================================

\author[Ash, Lou, Song]{%
    \texorpdfstring{%
    \begin{columns}
        \column{.3333\linewidth}
        \centering
        Elliot Ash \\\vspace{.5em} \small ETH Zurich
        \column{.3333\linewidth}
        \centering
        Yichuan Lou \\\vspace{.5em} \small UTokyo
        \column{.3333\linewidth}
        \centering
        Lena Song \\\vspace{.5em} \small UIUC
    \end{columns}}{Elliot Ash, Yichuan Lou, Lena Song}
}

\title[AI and Communication]{AI and Communication}
%\subtitle{Subtitle}
\date[Conference Name]{September 2025}




%====================================================
%========== BEGINNING OF DOCUMENT ===================
%====================================================
\begin{document}

\begin{frame}[plain]
    \titlepage
\end{frame}




\begin{frame}{Theoretical Framework: Overview}

\begin{itemize}\setlength{\itemsep}{1em}
    \item Framework based on Benabou and Tirole (2006) to model social signaling with image concerns in the presence of AI senders/receivers
    \item Two novel ingredients: outside option and privacy concerns
    \item Human sender makes a binary participation decision 
    \item Human sender makes his decision based on
    \begin{enumerate}
        \item intrinsic motivation
        \item participation cost
        \item image concerns
        \item privacy concerns
    \end{enumerate}
\end{itemize}
    
\end{frame}



\begin{frame}{Composition of Senders}

\begin{itemize}
    \item Three identity groups: humans, Gen-AI bots, and rule-based bots
    \item Denote by $m_h,m_g,m_r$ the number of different groups of senders 
\end{itemize}
    
\end{frame}



\begin{frame}{Timing}

\begin{enumerate}\setlength{\itemsep}{1em}
    \item A sender is matched with a receiver
    \item The sender's identity (human or bot) and type ($v\in [\underline{v},\overline{v}]$ if human) is initially unknown to the receiver
    \item The human sender chooses whether to participate in the signaling game
    \item If the human sender chooses to not participate, he obtains outside option with payoff $u_0$
    \item If the human sender chooses to participate, he incurs a fixed cost $c$
    \begin{itemize}
        \item the belief about the sender's identity and type is updated
        \item the human sender's payoff is then realized
    \end{itemize}
\end{enumerate}
    
\end{frame}


\begin{frame}{Human Sender}
If a human sender chooses to participate, he obtains
\begin{align*}
    U = v - c + \prob(\text{human}|\text{participate})\Big[\underbrace{\mu E(v|\text{participate})}_{\text{image concerns}} - \underbrace{\gamma}_{{\text{privacy concern}}}\Big]
\end{align*}
\vspace{1em}

\begin{itemize}
    \item $\prob(\text{human}|\text{participate})$ is the posterior probability that the sender is human
    \item $E(v|\text{participate})$ is the posterior expectation of the human sender's type $v$ 
    \item $\mu$ captures the strength of the human sender's image concerns
    \item $\gamma$ captures the human sender's concern about privacy leakage, i.e., the risk of revealing his identity (human v.s. bot), which is independent of the belief about $v$
\end{itemize}
\end{frame}


\begin{frame}{Bot Senders}
\begin{itemize}\setlength{\itemsep}{1em}
    \item A bot sender, whether Gen-AI or rule-based, has no private types or preferences
    \item Each bot sender follows a fixed action pattern
\end{itemize}
\vspace{1em}

\begin{exampleblock}{Assumption 1}
    A rule-based bot sender always chooses a detectable action $a_r$ that is neither 'participate' nor 'not participate'.
\end{exampleblock}

\begin{exampleblock}{Assumption 2}
    A Gen-AI sender always randomizes between 'participate' and a detectable action $a_g$: it plays 'participate' w.p. $\alpha$ and $a_g$ w.p. $(1-\alpha)$.
\end{exampleblock}
\end{frame}

\begin{frame}{Equilibrium with cutoff $v^*$}

We focus on threshold equilibria characterized by a cutoff $v^*$. In such an equilibrium, the human sender chooses
\begin{itemize}
    \item 'not participate' if $v\in [\underline{v},v^*)$
    \item 'participate' if $v\in (v^*,\overline{v}]$
\end{itemize}

\end{frame}

\begin{frame}{Human Sender Decision Making}

Let $v$ be drawn from a cdf $F\sim \Delta([\underline{v},\bar{v}])$
\vspace{1em}

Let $\rho\equiv\frac{m_g}{m_h+m_g}$ the relative probability of Gen-AI bots versus human senders

\vspace{1em}

Let $T(v^*)\equiv \mu E(v|v\geq v^*) - \gamma$

\vspace{1em}

Let $\pi(v^*)\equiv \frac{(1-\rho)(1-F(v^*)}{\rho\alpha + (1-\rho)(1-F(v^*))}$ be the probability that the sender is human conditional on participation

\vspace{1em}
Given any candidate cutoff $v^*$, a human sender with type $v$ participates iff
\begin{align*}
    v - c + \pi(v^*)T(v^*)\geq u_0
\end{align*}
    
\end{frame}

\begin{frame}{Equilibrium}

Let $\Phi(v^*)\equiv v^* - c + \pi(v^*)T(v^*)$
\vspace{1em}

The threshold equilibrium with cutoff $v^*$ satisfies
\begin{align*}
    \Phi(v^*) = u_0
\end{align*}

\begin{exampleblock}{Claim 1}
    In any stable equilibrium, $\Phi$ is increasing in $v^*$, i.e., $\frac{\partial \Phi(v^*)}{\partial v^*}>0$. 
\end{exampleblock}
\end{frame}

\begin{frame}{Some Calculations}

\begin{align*}
    \frac{\partial \Phi(v^*)}{\partial \rho} &= \frac{-\alpha(1-F(v^*))}{[\rho\alpha+(1-\rho)(1-F(v^*))]^2}T(v^*)\\[10pt]
    \frac{\partial \Phi(v^*)}{\partial \alpha} &= \frac{-\rho(1-\rho)(1-F(v^*))}{[\rho\alpha+(1-\rho)(1-F(v^*))]^2}T(v^*)\\[10pt]
    \frac{\partial \Phi(v^*)}{\partial \gamma} &= -\pi(v^*)<0
\end{align*}
\end{frame}

\begin{frame}{Interesting Case}
\begin{exampleblock}{Assumption 3}
    We focus on the interesting case where there exists an interior $v_c$ such that $T(v^*)<0$ for $v^*<v_c$ and $T(v^*)>0$ for $v^*>v_c$.
\end{exampleblock}
\vspace{2em}

\begin{itemize}
    \item when $T(v^*) = \mu E(v|v\geq v^*)-\gamma<0$, the privacy effect dominates the image effect
    \item when $T(v^*) = \mu E(v|v\geq v^*)-\gamma>0$, the image effect dominates the privacy effect
\end{itemize}
\end{frame}

\begin{frame}{Predictions}

\begin{itemize}
    \item when $T<0$, privacy dominates, then $v^*$ decreases in $\rho$
    \item when $T>0$, image dominates, then $v^*$ increases in $\rho$
    \item $v^*$ increases in $\gamma$
\end{itemize}





\end{frame}


\end{document}